\sscp{Sugar glider}{C-03-04}
\citet{Nivens2017} reconstructs Proto-\ili{Aru} **la[dŋ]i `sugarglider \it{Petaurus breviceps}'.
The \ili{Malay} gloss in \citet{Nivens2017} is ‘sejenis kuskus hitam
yang lincah meloncat di atas pohon’ = ‘kind of black cuscus which
leaps with agility between tree tops’.

\noindent{\wg\st{6pt}\tblr{llllll}
\lsptoprule
%Language	&	ISO	&	Term	&	Gloss	&	Etymon	&	Source	\\	\midrule
%\endfirsthead
Language	&	ISO	&	Term	&	Gloss	&	Etymon	&	Source	\\	\midrule
%\endhead
\ili{Ujir}	&	ugj	&	ladi	&	sugar glider	&	**la[dŋ]i	&	\cite{Nivens2017}	\\
\ili{Dobel}	&	kvo	&	laŋi	&	sugar glider	&	**la[dŋ]i	&	\cite{Nivens2017}	\\
NW. \ili{Tarangan}	&	txn	&	lan	&	sugar glider	&	**la[dŋ]i	&	\cite{Nivens2017}	\\
SW. \ili{Tarangan}	&	txn	&	lan	&	sugar glider	&	**la[dŋ]i	&	\cite{Nivens2017}	\\
\lspbottomrule
\end{tabular}}

%\noindent{\wg\st{2pt}\begin{tabularx}{\linewidth}
%{lllll>{\raggedright\arraybackslash}X}
%\lsptoprule
%%Language	&	ISO	&	Term	&	Gloss	&	Etymon	&	Source	\\	\midrule
%%\endfirsthead
%Language	&	ISO	&	Term	&	Gloss	&	Etymon	&	Source	\\	\midrule
%\endhead
%
%\lspbottomrule
%\end{tabularx}}

%\noindent{\wg\st{2pt}\begin{xltabular}{\linewidth}
%{lllll>{\raggedright\arraybackslash}X}
%\lsptoprule
%%Language	&	ISO	&	Term	&	Gloss	&	Etymon	&	Source	\\	\midrule
%%\endfirsthead
%Language	&	ISO	&	Term	&	Gloss	&	Etymon	&	Source	\\	\midrule
%\endhead
%
%\lspbottomrule
%\end{xltabular}}